% Cayley tomography and narrowing-orbit geometry for the resonance-window shadow packet.

\begin{corollary}[Cayley 断层测绘：由影子谱包直接读出 Breit--Wigner 的中心—宽度参数，并出现闭式常数 \texorpdfstring{$-\sqrt{15}/5$}{-sqrt15/5}]\label{cor:pom-resonance-cayley-breitwigner-centerwidth}
\leanverified{paper\_pom\_resonance\_cayley\_breitwigner\_centerwidth}
在注 \ref{rem:pom-resonance-poisson-lorentz-prelimit} 的设定下，对任意盘内点 \(a=\eta e^{\mathrm{i}\theta}\in\mathbb{D}\) 定义 Cayley 半平面像
\[
w:=\frac{1+a}{1-a}=\delta+\mathrm{i}\gamma\qquad(\delta>0).
\]
则有显式公式
\[
\boxed{\;
\delta=\frac{1-\eta^{2}}{1-2\eta\cos\theta+\eta^{2}},
\qquad
\gamma=\frac{2\eta\sin\theta}{1-2\eta\cos\theta+\eta^{2}}\; }.
\]
在内函数相位导数的 Poisson（Breit--Wigner）叠加律中，\(w\) 对应一项 Lorentz 核
\(
\delta\big/((t-\gamma)^2+\delta^2)
\)
的“峰位 \(\gamma\) 与半宽 \(\delta\)”参数。
\par
将此应用于共振窗影子谱包，取
\(
a_q:=\eta_q e^{\mathrm{i}\theta_q}=\mu_q/r_{q-2}
\)
并记 \(w_q=\delta_q+\mathrm{i}\gamma_q\)。
若进一步满足 \(\eta_q\to 1\) 且 \(\theta_q\to \theta^\ast:=\arg(z_\ast)\)（其中 \(z_\ast=(-1-\mathrm{i}\sqrt{15})/4\)，故 \(\cos\theta^\ast=-\tfrac14\)），则 \(w_q\to \mathrm{i}\gamma^\ast\) 且
\[
\boxed{\;
\gamma^\ast=\cot\!\Bigl(\frac{\theta^\ast}{2}\Bigr)=-\frac{\sqrt{15}}{5},
\qquad
\delta_q\sim \frac{4}{5}\,(1-\eta_q)\; }.
\]
\end{corollary}
\begin{proof}
由
\(
\frac{1+a}{1-a}=\frac{(1+a)(1-\overline a)}{|1-a|^2}
\)
得
\(
\delta=\Re(w)=\frac{1-|a|^2}{|1-a|^2}
\)
与
\(
\gamma=\Im(w)=\frac{2\Im(a)}{|1-a|^2}
\)；
代入 \(a=\eta e^{\mathrm{i}\theta}\) 并写 \(|1-a|^2=1-2\eta\cos\theta+\eta^2\) 即得两条显式公式。
\par
若 \(\eta_q\to 1\) 且 \(\theta_q\to\theta^\ast\)，则 \(a_q\to e^{\mathrm{i}\theta^\ast}\)，从而
\(
w_q\to\frac{1+e^{\mathrm{i}\theta^\ast}}{1-e^{\mathrm{i}\theta^\ast}}
=\mathrm{i}\cot(\theta^\ast/2)
\)，即 \(w_q\to \mathrm{i}\gamma^\ast\)。
由 \(\cos\theta^\ast=-\tfrac14\)、\(\sin\theta^\ast=-\tfrac{\sqrt{15}}{4}\) 得
\[
\gamma^\ast=\cot\!\Bigl(\frac{\theta^\ast}{2}\Bigr)=\frac{1+\cos\theta^\ast}{\sin\theta^\ast}
=\frac{3/4}{-\sqrt{15}/4}=-\frac{\sqrt{15}}{5}.
\]
最后，\(\delta=\frac{1-\eta^2}{1-2\eta\cos\theta+\eta^2}\) 在 \((\eta,\theta)\to(1,\theta^\ast)\) 下满足
\[
\delta\sim \frac{(1-\eta)(1+\eta)}{2-2\cos\theta^\ast}
\sim \frac{2(1-\eta)}{2(1-\cos\theta^\ast)}
=\frac{1-\eta}{1-\cos\theta^\ast}
=\frac{4}{5}(1-\eta),
\]
其中用到 \(1-\cos\theta^\ast=1+\tfrac14=\tfrac54\)。证毕。
\end{proof}

\begin{corollary}[共振窗的几何本质：固定峰位的窄化轨道——“能量几乎不动、宽度急剧收缩”（可复算）]\label{cor:pom-resonance-cayley-narrowing-orbit}
\leanverified{paper\_pom\_resonance\_cayley\_narrowing\_orbit}
对共振窗 \(q=9,\dots,17\)，取表 \ref{tab:fold_collision_shadow_spectral_packet_q9_17} 的两步模量与相位
\(
\eta_q=|\mu_q|/r_{q-2}
\)、\(\theta_q=\arg(\mu_q)
\)，并按推论 \ref{cor:pom-resonance-cayley-breitwigner-centerwidth} 定义
\[
w_q=\frac{1+a_q}{1-a_q}=\delta_q+\mathrm{i}\gamma_q,
\qquad a_q:=\eta_q e^{\mathrm{i}\theta_q}.
\]
则由表中数据直接代入可得：
\begin{enumerate}
  \item \textbf{峰位几乎固定：}\(\gamma_q\) 被压缩在长度
  \[
  \max_{9\le q\le 17}\gamma_q-\min_{9\le q\le 17}\gamma_q\ <\ 1.1\times 10^{-2}
  \]
  的短区间内（数值上 \(\gamma_q\in[-0.7711,-0.7608]\)）。
  \item \textbf{半宽单调窄化：}\(\delta_q\) 随 \(q\) 单调下降，并在窗端达到
  \[
  \delta_{17}\approx 3.16\times 10^{-3}=O(10^{-3}).
  \]
\end{enumerate}
因此影子谱包在 Cayley 半平面中形成一条逼近边界的窄化轨道：其峰位（中心）几乎不动，而半宽随 \(1-\eta_q\) 快速闭合。
\end{corollary}
\begin{proof}
由推论 \ref{cor:pom-resonance-cayley-breitwigner-centerwidth} 的显式公式
\[
\delta_q=\frac{1-\eta_q^{2}}{1-2\eta_q\cos\theta_q+\eta_q^{2}},
\qquad
\gamma_q=\frac{2\eta_q\sin\theta_q}{1-2\eta_q\cos\theta_q+\eta_q^{2}},
\]
将表 \ref{tab:fold_collision_shadow_spectral_packet_q9_17} 的 \((\eta_q,\theta_q)\)（四舍五入读数）逐项代入即可得到所述区间与单调性结论。证毕。
\end{proof}

\endinput
